\documentclass[11pt,a4paper]{article}

% ============================================================================
% DRAFT WATERMARK - PROMINENT LABELING
% ============================================================================
\usepackage{draftwatermark}
\SetWatermarkText{WORKING DRAFT}
\SetWatermarkScale{0.5}
\SetWatermarkColor[gray]{0.85}

\usepackage[margin=1in]{geometry}
\usepackage{amsmath,amssymb,amsthm}
\usepackage{hyperref}
\usepackage{booktabs}
\usepackage{graphicx}
\usepackage{xcolor}
\usepackage{listings}
\usepackage{algorithm}
\usepackage{algpseudocode}

% Draft notice box
\usepackage{tcolorbox}
\newtcolorbox{draftnotice}{
  colback=red!5!white,
  colframe=red!75!black,
  fonttitle=\bfseries,
  title=WORKING DRAFT
}

% Theorem environments
\newtheorem{theorem}{Theorem}[section]
\newtheorem{lemma}[theorem]{Lemma}
\newtheorem{definition}[theorem]{Definition}
\newtheorem{conjecture}[theorem]{Conjecture}
\newtheorem{corollary}[theorem]{Corollary}
\newtheorem{example}[theorem]{Example}

% Code listings
\lstset{
  basicstyle=\ttfamily\small,
  breaklines=true,
  frame=single,
  language=Python
}

\title{Quality-Guided Descent: Deterministic Gates for\\Stochastic LLM Agents}

\author{
  Andrew Templeton\\
  \texttt{github.com/andrew-templeton/quality-gate-sgd}
}

\date{Working Draft -- \today}

\begin{document}

\maketitle

% ============================================================================
% PROMINENT DRAFT NOTICE
% ============================================================================
\begin{draftnotice}
This document is a \textbf{working draft} of preliminary theoretical work.
\begin{itemize}
  \item Claims marked \textbf{[NOVEL]} await empirical validation
  \item Citations marked \textbf{[PENDING]} are being compiled
  \item This work has \textbf{not been peer-reviewed}
\end{itemize}
See the companion repository for the latest version and claim status:\\
\url{https://github.com/andrew-templeton/quality-gate-sgd}
\end{draftnotice}

\begin{abstract}
We present a theoretical framework for understanding how deterministic quality gates create gradient descent-like optimization behavior when LLM agents iterate against them. We formalize the required properties of quality functions (quantitative measurement, determinism, local continuity), prove a convergence theorem under a ``biased proposer'' assumption, and introduce \emph{discrete differentiability}--a method for computing target-space gradients over enumerable code locations. This enables optimization guidance that specifies not just \emph{which} metric to improve, but \emph{where} in the codebase to make changes. We classify common software quality metrics by their suitability for this framework and propose empirically testable hypotheses about convergence rates under different feedback granularities.

\medskip
\noindent\textbf{Status:} Working draft. Novel claims require experimental validation.
\end{abstract}

\tableofcontents
\newpage

% ============================================================================
\section{Introduction}
% ============================================================================

Large Language Models (LLMs) have demonstrated remarkable capability in code generation and modification tasks. However, their outputs are inherently stochastic--the same prompt may yield different code across invocations. This stochasticity poses challenges for software engineering workflows that require deterministic, reproducible quality outcomes.

We propose that the path to deterministic results from stochastic agents lies not in constraining the agent, but in making the \emph{exit gate} deterministic. When an LLM agent iterates against a deterministic quality function, rejecting outputs that fail to meet quality thresholds, the system exhibits behavior analogous to stochastic gradient descent (SGD):

\begin{itemize}
  \item The quality function provides the \emph{descent direction} (what to improve)
  \item The LLM's stochasticity provides \emph{exploration} (diverse fix attempts)
  \item Gate rejection creates \emph{optimization pressure} (improve or fail)
\end{itemize}

This paper formalizes this intuition, establishing conditions under which convergence is guaranteed and introducing computational methods for fine-grained optimization guidance.

% ============================================================================
\section{Quality Geometry}
\label{sec:geometry}
% ============================================================================

\subsection{The Quality Space}

\begin{definition}[Quality Function]
Let $\mathcal{C}$ be the space of all syntactically valid code states for a project. A \textbf{quality function} is a mapping:
\[
Q: \mathcal{C} \to \mathbb{R}^n
\]
where each dimension $q_i$ represents a measurable quality attribute (e.g., test coverage, bug count, cyclomatic complexity).
\end{definition}

\subsection{Required Properties}

For $Q$ to support gradient-like optimization, it should satisfy three properties:

\begin{definition}[Quantitative Measurement]
$Q(c) \in \mathbb{R}^n$ for all $c \in \mathcal{C}$. Each dimension must produce a real number with a clear optimization direction.
\end{definition}

\begin{definition}[Determinism]
$Q(c_1) = Q(c_2)$ whenever $c_1 = c_2$. The same code state must always produce the same quality vector.
\end{definition}

\begin{definition}[Local Continuity]
For all $\varepsilon > 0$, there exists $\delta > 0$ such that:
\[
d(c_1, c_2) < \delta \implies \|Q(c_1) - Q(c_2)\| < \varepsilon
\]
where $d$ is a distance metric on code space (e.g., edit distance).
\end{definition}

\subsection{The Descent Condition}

When these properties hold, a quality space admits \emph{descent}: there exist transformation sequences that monotonically improve $Q$.

\begin{theorem}[Descent Existence, Informal]
If $Q$ satisfies Properties 1-3, then for most code states $c$ where $Q(c)$ is not at a local minimum, there exists a neighborhood $N(c)$ containing states $c'$ with $Q(c') \prec Q(c)$ (where $\prec$ denotes improvement).
\end{theorem}

% ============================================================================
\section{Convergence Theory}
\label{sec:convergence}
% ============================================================================

\subsection{The Agent Model}

\begin{definition}[Feedback Function]
Given current measurements $Q(c)$, the feedback function $F$ identifies constraint violations:
\[
F: \mathcal{M} \to \text{Feedback}
\]
where $\mathcal{M} \subseteq \mathbb{R}^n$ is the measurement space. This function is deterministic.
\end{definition}

\begin{definition}[LLM Agent as Stochastic Proposer]
An LLM agent $A$ takes current code and feedback, proposing a modification:
\[
A: \mathcal{C} \times \text{Feedback} \to \text{Distribution}(\mathcal{C})
\]
Given $(c, F(Q(c)))$, the agent samples $c' \sim A(c, F(Q(c)))$.
\end{definition}

\subsection{The Biased Proposer Property}

\begin{definition}[Biased Proposer]
\label{def:biased-proposer}
An agent $A$ is a \textbf{biased proposer} with parameter $p > 0.5$ if:
\[
P(Q(c') \succ Q(c) \mid c' \sim A(c, F(Q(c)))) \geq p
\]
That is, given feedback about failing metrics, the agent proposes improvements with probability at least $p$.
\end{definition}

\textbf{Claim [NOVEL]:} Modern LLMs (GPT-4, Claude) are biased proposers for code quality tasks. This is empirically testable.

\subsection{The Convergence Theorem}

\begin{theorem}[Finite Expected Convergence]
\label{thm:convergence}
Let:
\begin{itemize}
  \item $Q: \mathcal{C} \to \mathcal{M}$ be a quality function with finite $|\mathcal{M}|$
  \item $T \subseteq \mathcal{M}$ be a non-empty target set (satisfying all constraints)
  \item $A$ be a biased proposer with parameter $p > 0.5$
  \item The quality ordering $\prec$ have no infinite descending chains in $\mathcal{M} \setminus T$
\end{itemize}

Then starting from any $c_0$ with $Q(c_0) \notin T$:
\[
\mathbb{E}[\tau] < \infty
\]
where $\tau = \inf\{t : Q(c_t) \in T\}$ is the stopping time.
\end{theorem}

\begin{proof}[Proof Sketch]
\begin{enumerate}
  \item $\mathcal{M}$ is finite and $\prec$ has no infinite descending chains, so there exists a maximum distance $d_{\max}$ from any $m \in \mathcal{M}$ to $T$.
  \item At each step, with probability $\geq p > 0.5$, we improve (move closer to $T$).
  \item This is a random walk with drift toward an absorbing set.
  \item By standard results on biased random walks [PENDING: Feller, Ross], absorption occurs in finite expected time.
  \item Expected proposals per accepted step is $1/p$ (geometric distribution).
  \item Therefore $\mathbb{E}[\tau] \leq d_{\max}/p < \infty$. \qedhere
\end{enumerate}
\end{proof}

\begin{corollary}[Convergence Rate]
If $d_{\max}$ is the maximum distance from any $m \in \mathcal{M} \setminus T$ to $T$, then:
\[
\mathbb{E}[\tau] \leq \frac{d_{\max}}{p}
\]
\end{corollary}

% ============================================================================
\section{Discrete Differentiability}
\label{sec:differentiability}
% ============================================================================

\subsection{The Problem: Gradient Degeneracy}

The standard quality gradient operates in \emph{dimension-space}:
\[
\nabla Q = \left[\frac{\partial Q}{\partial q_1}, \frac{\partial Q}{\partial q_2}, \ldots, \frac{\partial Q}{\partial q_n}\right]
\]

This tells us \emph{which} dimension to improve but not \emph{where} in the codebase to make changes. In discrete spaces, this degenerates to priority ordering over dimensions.

\subsection{The Insight: Location Data Exists}

Quality measurement tools already extract precise location information:

\begin{center}
\begin{tabular}{lll}
\toprule
\textbf{Source} & \textbf{Location Data} & \textbf{Current Usage} \\
\midrule
coverage-final.json & file, line, branch, function & Aggregated to \% \\
TypeScript errors & file:line:column & Counted only \\
ESLint issues & file:line:column:ruleId & Counted only \\
SonarQube issues & file:line:rule:severity & Counted only \\
\bottomrule
\end{tabular}
\end{center}

\textbf{Key realization:} We have the data to compute gradients over code-space. We discard it during aggregation.

\subsection{Target-Space Gradients}

\begin{definition}[Optimization Target]
Let $\mathcal{T} = \{t_1, t_2, \ldots, t_m\}$ be the set of enumerable optimization targets induced by the addressing scheme, where each target can be represented as a tuple:
\[
t = (\text{address\_id}, \text{span}?, \text{issue\_cluster})
\]
\end{definition}

\begin{definition}[Target-Space Gradient]
The \textbf{target-space gradient} is:
\[
\nabla_t Q = \left[\frac{\partial Q}{\partial t_1}, \frac{\partial Q}{\partial t_2}, \ldots, \frac{\partial Q}{\partial t_m}\right]
\]
where $\partial Q / \partial t_i$ represents the expected change in fitness if all issues at target $t_i$ are resolved.
\end{definition}

\subsection{Addressing Schemes and Consistency}

Target-space gradients only become actionable when all issue sources map into a shared address space. We call this the \emph{addressing scheme}. This choice is part of the topology: it specifies the coordinate system over which discrete gradients are computed.

\begin{definition}[Addressing Scheme]
An addressing scheme is a deterministic mapping
\[
A: \mathcal{C} \to \mathcal{A}
\]
that yields an address space $\mathcal{A} = (V, E, \mu)$ where $V$ is the set of addressable units, $E$ is an optional adjacency relation, and $\mu$ assigns a size measure used for normalization. We take $\mathcal{T} = V$ as the target set for target-space gradients.
\end{definition}

\begin{definition}[Quality Coordinate Space]
An address space $\mathcal{A}$ is a \emph{quality coordinate space} if it is:
\begin{itemize}
  \item \textbf{Deterministic}: the same code state yields the same addresses
  \item \textbf{Stable}: small edits preserve address identity for unchanged units
  \item \textbf{Shared}: every metric axis can map its issues into $V$
  \item \textbf{Normalizable}: each address has a size measure $\mu$ for density
  \item \textbf{Graphable (optional)}: $E$ enables weighting and propagation
\end{itemize}
\end{definition}

In practice, a scheme may fall back to coarser addresses (e.g., file:line) when a precise address is unavailable.

For TypeScript, we use the compiler symbol graph as the address space: symbols are nodes, edges (calls, imports, extends) provide adjacency, and symbol span gives $\mu$. File-and-line addressing is simpler but tends to be fragile under refactors, so we use it only as a fallback.

The same requirement applies to non-code domains: a proof may be addressed by a claim graph, and an essay by paragraph/sentence/word indices or a topic graph, but the scheme must be consistent across all metrics.

\subsection{Computing $\Delta Q$}

For a target $t$ with issues affecting dimensions $D = \{d_1, d_2, \ldots, d_k\}$:

\begin{equation}
\Delta Q(t) = \sum_{j \in D} w_j \cdot \delta_j(t)
\label{eq:delta-q}
\end{equation}

where:
\begin{itemize}
  \item $w_j$ is the weight for dimension $j$
  \item $\delta_j(t)$ is the normalized impact on dimension $j$ from fixing issues at target $t$
\end{itemize}

\subsection{Cross-Dimension Correlation}

A target affecting multiple dimensions simultaneously is more valuable than one addressing a single dimension.

\begin{example}
Consider two targets with similar coverage gaps:
\begin{center}
\begin{tabular}{lcccc}
\toprule
\textbf{Target} & \textbf{Coverage $\Delta Q$} & \textbf{TS Errors $\Delta Q$} & \textbf{Smells $\Delta Q$} & \textbf{Total $\Delta Q$} \\
\midrule
\texttt{processPayment()} & +0.35 & +0.25 & +0.10 & \textbf{0.70} \\
\texttt{formatDate()} & +0.07 & 0 & 0 & \textbf{0.07} \\
\bottomrule
\end{tabular}
\end{center}

The first target is \textbf{10$\times$ more valuable} because it addresses three dimensions. Dimension-level gradients cannot see this correlation.
\end{example}

\subsection{Granularity and Convergence}

\begin{conjecture}[Granularity-Convergence Relationship]
\label{conj:granularity}
Let $\tau_{\text{dim}}$, $\tau_{\text{file}}$, $\tau_{\text{symbol}}$ be expected iterations to convergence for dimension-level, file-level, and symbol-level feedback respectively. Then:
\[
\tau_{\text{symbol}} < \tau_{\text{file}} < \tau_{\text{dim}}
\]
\end{conjecture}

\textbf{Rationale:} Finer-grained feedback reduces agent search space:
\begin{itemize}
  \item ``Improve coverage'' $\to$ agent searches entire codebase
  \item ``Fix \texttt{src/payment.ts}'' $\to$ agent focuses on one file
  \item ``Fix \texttt{processPayment()}'' $\to$ agent focuses on one function
\end{itemize}

This conjecture is empirically testable.

% ============================================================================
\section{Metric Topology}
\label{sec:topology}
% ============================================================================

\subsection{Metric Classification}

We classify metrics by their suitability for the SGD framework:

\begin{center}
\begin{tabular}{lcccl}
\toprule
\textbf{Metric} & \textbf{Quant.} & \textbf{Determ.} & \textbf{Contin.} & \textbf{Grade} \\
\midrule
coverage.lines & $\checkmark$ & $\checkmark$ & $\checkmark\checkmark$ & A \\
coverage.branches & $\checkmark$ & $\checkmark$ & $\checkmark$ & B+ \\
duplications \% & $\checkmark$ & $\checkmark$ & $\checkmark$ & B+ \\
sonarqube.codeSmells & $\checkmark$ & $\checkmark$ & $\triangle$ & B- \\
sonarqube.bugs & $\checkmark$ & $\checkmark$ & $\times$ & C- \\
sonarqube.blocker & $\checkmark$ & $\checkmark$ & $\times$ & D- \\
typescript.errors & $\checkmark$ & $\triangle^*$ & $\times$ & D \\
\bottomrule
\end{tabular}
\end{center}

$^*$TypeScript errors can be non-deterministic with incremental compilation.

\subsection{Design Principles}

Based on this analysis, we recommend:

\begin{enumerate}
  \item \textbf{Smooth objectives}: Use percentage-based metrics (coverage) as optimization targets
  \item \textbf{Discrete constraints}: Use count metrics (blockers, criticals) as ceilings, not objectives
  \item \textbf{Normalization}: Transform counts to per-kSLOC densities for improved continuity
  \item \textbf{Weighting}: Use dependency analysis to focus effort on high-impact code
\end{enumerate}

% ============================================================================
\section{Related Work}
\label{sec:related}
% ============================================================================

\textbf{[PENDING: Citations to be compiled]}

\begin{itemize}
  \item \textbf{RLHF and LLM Alignment}: InstructGPT, Constitutional AI
  \item \textbf{LLM Code Generation}: Codex, HumanEval benchmark
  \item \textbf{Random Walk Convergence}: Feller, Ross
  \item \textbf{Software Quality Metrics}: ISO/IEC 25010, SonarQube validation studies
  \item \textbf{Multi-objective Optimization}: Pareto optimality literature
\end{itemize}

% ============================================================================
\section{Empirical Validation Plan}
\label{sec:validation}
% ============================================================================

\subsection{Research Questions}

\begin{description}
  \item[RQ1: Convergence] Do LLM agents with quality gates converge faster than agents without?
  \item[RQ2: Success Rate] Do quality gates improve task completion rate?
  \item[RQ3: Quality Ceiling] Do quality gates improve final code quality beyond ``just passing''?
  \item[RQ4: Topology Design] Which metric combinations create the smoothest descent?
  \item[RQ5: Target Granularity] Does file-level or symbol-level targeting improve convergence?
  \item[RQ6: Cross-Dimension Value] Do multi-dimension targets converge faster?
  \item[RQ7: Addressing Fitness] Does addressing fitness predict convergence speed and success rate?
\end{description}

\subsection{Biased Proposer Validation}

\textbf{Protocol:}
\begin{enumerate}
  \item Collect $N$ code states with known quality issues
  \item For each state $c_i$, generate feedback $F(Q(c_i))$
  \item Ask LLM to propose fix $c'_i$
  \item Measure: did $Q(c'_i) \succ Q(c_i)$?
  \item Compute $\hat{p} = (\text{\# improvements}) / N$
\end{enumerate}

\textbf{Hypothesis:} $\hat{p} > 0.5$ with statistical significance.

\textbf{Required sample size:} $N \approx 100$ for 95\% confidence.

% ============================================================================
\section{Conclusion}
\label{sec:conclusion}
% ============================================================================

We have presented a theoretical framework for understanding how deterministic quality gates create optimization behavior in stochastic LLM agents. Key contributions:

\begin{enumerate}
  \item \textbf{Quality Geometry}: Formalization of the required properties (quantitative, deterministic, continuous)
  \item \textbf{Convergence Theorem}: Proof of finite expected convergence under the biased proposer assumption
  \item \textbf{Discrete Differentiability}: Method for computing target-space gradients over code locations
  \item \textbf{Metric Topology}: Classification of common metrics by SGD suitability
\end{enumerate}

The framework provides both theoretical grounding and practical guidance for designing quality-gated LLM coding workflows. All novel claims are explicitly marked for empirical validation.

% ============================================================================
\section*{Acknowledgments}
% ============================================================================

This is a working draft. The author welcomes feedback and collaboration on empirical validation.

% ============================================================================
\appendix
\section{Claim Inventory}
\label{app:claims}
% ============================================================================

\begin{center}
\begin{tabular}{lll}
\toprule
\textbf{Claim} & \textbf{Type} & \textbf{Status} \\
\midrule
Quality functions map code to finite $\mathcal{M}$ & MATH & Self-supporting \\
Finite granularity induces equivalence classes & MATH & Self-supporting \\
LLMs are biased proposers ($p > 0.5$) & NOVEL & Requires validation \\
Convergence in finite expected time & NOVEL & Theorem~\ref{thm:convergence} \\
Target-space gradients improve guidance & NOVEL & Requires validation \\
Finer granularity $\to$ faster convergence & NOVEL & Conjecture~\ref{conj:granularity} \\
Random walks with drift converge & CITED & PENDING \\
RLHF improves LLM behavior & CITED & PENDING \\
\bottomrule
\end{tabular}
\end{center}

\end{document}
